% Solution: Mas-Colell, Whinston & Green (1995) Exercise 9.B.5 (E. Maskin). Input by \solitem.
% New in the 30-item list. Check: check_B_L04.py. Compared with the MWG solutions manual.
\begin{exsol}{MWG Exercise 9.B.5}{302}{Advanced}{5}
\asked For a finite two-player game played simultaneously or with player~1 moving first and
observed: (a), (b) count strategies; (c) multiple SPNE imply a payoff tie; (d) without ties for
player~2, moving first guarantees player~1 at least any pure Nash payoff of the simultaneous game,
in every SPNE (and whether in every Nash equilibrium); (e) examples where (d) fails with ties for
player~2 or with a mixed Nash equilibrium.

\method Strategies as complete plans (MWG \S7.D; Lecture~3, slide~25); SPNE and backward induction in
finite games of perfect information (MWG \S9.B, Prop.~9.B.1 (Zermelo), p.~272, and Prop.~9.B.2;
Lecture~3, slides~30 and~35; Lecture~4, slide~17). The sequential game in (b) has perfect information,
so its SPNE are the backward-induction profiles; with no ties they are unique. Non-credible threats in
Nash equilibrium: Lecture~4, slides~8 and~13.

\solution
\textbf{Notation.} As in the statement, $m_i$ is both the number of actions and (from (c) on) a typical
action; write $|M_i|$ for the number to avoid confusion. In the sequential game a strategy of
player~2 is a function $s_2: M_1 \to M_2$.

\textbf{(a) Simultaneous moves.} Each player has one information set, so player~$i$ has $|M_i| = m_i$
pure strategies (and the mixed strategies $\Delta(M_i)$, a simplex of dimension $m_i - 1$).

\textbf{(b) Player 1 first, observed.} Player~1 still has one information set: $m_1$ pure strategies.
Player~2 has one information set for each of player~1's moves, each with $m_2$ actions, so
$m_2^{\,m_1}$ pure strategies (the functions $M_1 \to M_2$; item ``Tadelis 7.1'').

\textbf{(c) Multiple SPNE imply a tie.} Suppose, to the contrary, that no player is indifferent between
any two distinct pairs: $\phi_i(m_1, m_2) \ne \phi_i(m_1', m_2')$ for $i = 1,2$ whenever
$(m_1,m_2) \ne (m_1',m_2')$.
\begin{itemize}
  \item \emph{Player~2's subgames.} After each $m_1$, the subgame is a one-player decision; any SPNE
  must play a best reply there. Since $\phi_2(m_1, \cdot)$ takes $m_2$ distinct values, the best reply
  $b(m_1)$ is unique and played with probability 1. So player~2's SPNE strategy is the unique
  function $b$.
  \item \emph{Player~1.} Given $b$, choosing $m_1$ yields $\phi_1(m_1, b(m_1))$. The pairs
  $(m_1, b(m_1))$ differ across $m_1$ (their first components differ), so by assumption these
  $m_1$ payoffs are distinct, and the maximizer $m_1^*$ is unique; any SPNE puts probability 1 on it.
\end{itemize}
So the SPNE is unique, $(m_1^*, b)$, contradicting multiplicity. Hence if there are multiple SPNE,
some player has a tie: (i) or (ii) holds for two distinct pairs. (In fact more is shown: either
player~2 has a tie \emph{within} some $\phi_2(m_1, \cdot)$, or player~1 has a tie between two pairs
with different $m_1$.) $\square$

\textbf{(d) Moving first does not hurt player 1 (in SPNE).} Assume (ii) never holds: player~2 has no
ties. Let $(\hat m_1, \hat m_2)$ be a pure Nash equilibrium of the simultaneous game with
$\pi_1 = \phi_1(\hat m_1, \hat m_2)$.
\begin{enumerate}[label=(\roman*)]
  \item In the simultaneous game $\hat m_2$ is a best reply to $\hat m_1$: $\phi_2(\hat m_1, \hat m_2)
  \ge \phi_2(\hat m_1, m_2)$ for all $m_2$. With no ties, it is the \emph{unique} best reply, so
  $b(\hat m_1) = \hat m_2$.
  \item In any SPNE of the sequential game, player~2 plays $b(m_1)$ after every $m_1$ (with
  probability 1, since best replies are unique). In particular she plays $\hat m_2$ after $\hat m_1$.
  \item An SPNE is a Nash equilibrium, so player~1's equilibrium payoff is at least what he gets by
  deviating to $\hat m_1$, which is $\phi_1(\hat m_1, b(\hat m_1)) = \phi_1(\hat m_1, \hat m_2) =
  \pi_1$.
\end{enumerate}
So in every SPNE player~1 gets at least $\pi_1$ (for every pure Nash equilibrium of (a), hence at least
the largest such $\pi_1$). $\square$

\emph{Not for every Nash equilibrium of (b).} A Nash equilibrium of the sequential game may use a
non-credible threat by player~2 after moves that player~1 does not make. Example (no ties for
player~2; rows player~1, columns player~2, both choose $O$ or $F$):
\begin{center}
\begin{tabular}{r|c|c|}
\multicolumn{1}{r}{} & \multicolumn{1}{c}{$O$} & \multicolumn{1}{c}{$F$}\\ \cline{2-3}
$O$ & $2,\ 1$ & $0,\ 0$ \\ \cline{2-3}
$F$ & $0,\ -1$ & $1,\ 2$ \\ \cline{2-3}
\end{tabular}
\end{center}
Player~2's payoffs $1, 0, -1, 2$ are all distinct. The simultaneous game has the pure equilibrium
$(O, O)$ with $\pi_1 = 2$ (and $(F,F)$ with 1). In the sequential game the unique SPNE is
$(O,\ s_2 = (O \text{ after } O, F \text{ after } F))$, payoff 2 to player~1, as (d) predicts. But
$(F,\ s_2 = (F, F))$ is also a Nash equilibrium: against ``always $F$'', player~1 gets $0$ from $O$ and
$1$ from $F$; against $F$, player~2's on-path choice $F$ is optimal ($2 > -1$). Player~1 gets
$1 < 2 = \pi_1$. The threat ``$F$ after $O$'' ($0 < 1$ for player~2) is not credible.

\textbf{(e) Two ways the conclusion of (d) fails.}

\emph{(e.1) Player~2 has ties.} Take the game of the MWG manual, with player~2's payoff depending
only on player~1's move:
\begin{center}
\begin{tabular}{r|c|c|}
\multicolumn{1}{r}{} & \multicolumn{1}{c}{$L$} & \multicolumn{1}{c}{$R$}\\ \cline{2-3}
$U$ & $0,\ 1$ & $1,\ 1$ \\ \cline{2-3}
$D$ & $0,\ 0$ & $1,\ 0$ \\ \cline{2-3}
\end{tabular}
\end{center}
$(U, R)$ is a pure Nash equilibrium of the simultaneous game with $\pi_1 = 1$ (player~1: $1 \ge 1$;
player~2: $1 \ge 1$). In the sequential game player~2 is indifferent at both nodes, so every reply is
sequentially rational, and $(U, (L \text{ after } U, L \text{ after } D))$ is an SPNE (player~1 gets 0
from either move). Player~1's SPNE payoff is $0 < 1 = \pi_1$.

\emph{(e.2) A mixed Nash equilibrium instead of a pure one.} Matching pennies:
\begin{center}
\begin{tabular}{r|c|c|}
\multicolumn{1}{r}{} & \multicolumn{1}{c}{$L$} & \multicolumn{1}{c}{$R$}\\ \cline{2-3}
$U$ & $1,\ 0$ & $0,\ 1$ \\ \cline{2-3}
$D$ & $0,\ 1$ & $1,\ 0$ \\ \cline{2-3}
\end{tabular}
\end{center}
Player~2 has no ties within either row. The simultaneous game has no pure equilibrium; its mixed
equilibrium $(\frac12,\frac12)$ gives player~1 $\frac12$. In the sequential game player~2 replies $R$
after $U$ and $L$ after $D$, so player~1 gets 0 whatever he does: the SPNE are $(U, (R, L))$ and
$(D, (R, L))$, both with payoff $0 < \frac12$. Moving first and being observed is a
\emph{disadvantage} when player~1 wants to be unpredictable.

\interp Moving first helps when it is a commitment the follower must respond to (Stackelberg, entry
deterrence); (d) makes this precise: if the follower's reply to each move is pinned down, the leader
can always reproduce any pure Nash outcome of the simultaneous game by playing his part of it, so he
cannot lose by moving first. The two failures in (e) mark the limits: if the follower is indifferent,
her tie-breaking can punish the leader at no cost to herself; and if the leader's simultaneous-move
advantage comes from being unpredictable (mixing), observation destroys it. The Nash-but-not-SPNE
example in (d) is the Lecture~4 point once more: threats off the path can make the leader look worse
off, but they are not credible.

\verify \texttt{check\_B\_L04.py} enumerates all pure SPNE (a strategy $s_2: M_1 \to M_2$ that best
replies in every subgame, and a best $m_1$) of 3000 random games with $m_1, m_2 \in \{2, 3\}$ and
integer payoffs in $\{0, \dots, 4\}$ (fixed seed). In each of the 1795 games with several SPNE a tie
(i) or (ii) is present; in each of the 166 games without ties for player~2, every SPNE gives player~1
at least the payoff of every pure Nash equilibrium of the simultaneous game. It confirms the
counterexample in (d) (unique SPNE pays 2, the Nash equilibrium $(F, (F,F))$ pays 1), (e.1) (SPNE
$(U, (L,L))$ pays 0 against $\pi_1 = 1$) and (e.2) (SPNE $(U,(R,L))$, $(D,(R,L))$ pay 0; mixed value
$\frac12$). Output line: \texttt{mwg-9-b-5: (c) on 1795 random games with multiple SPNE, (d) on 166
games; examples (d), (e1), (e2) confirmed}. \emph{Comparison with the MWG solutions manual}
(Exercise 9.B.5, pp.~9-2 to 9-4; label OCR'd as ``9.8.5''): (a) $m_i$ pure strategies; (b) $m_1$ and
$m_2^{m_1}$; (c) proof by contradiction, as here; (d) the same argument, and ``low payoffs for player~1
could be sustained in a NE with incredible threats by player~2''; (e) the same two examples (the
indifferent follower of Fig.~9.B.5(e.1) and matching pennies, Fig.~9.B.5(e.2)). \textbf{Agreement on
every part.} One slip in the manual (or its OCR): the SPNE of the pennies game are written
``$(U,(R,L))$ and $(L,(R,L))$''; the second must be $(D, (R,L))$, since $L$ is not a move of player~1.

\mistakes
\begin{itemize}
  \item Giving player~2 $m_1 m_2$ strategies in (b) (pairs) instead of $m_2^{m_1}$ (functions).
  \item In (d), forgetting that the argument needs player~2's best reply to $\hat m_1$ to be
  \emph{unique}: with a tie she might answer $\hat m_1$ differently in the sequential game, which is
  exactly example (e.1).
  \item Claiming that moving first is always an advantage: (e.2) shows it can be strictly worse.
\end{itemize}

\answer{(a) $m_1$ and $m_2$. (b) $m_1$ for player~1, $m_2^{\,m_1}$ for player~2. (c) Without ties
player~2's reply to each move is unique and then player~1's best move is unique, so the SPNE is unique;
multiple SPNE therefore require (i) or (ii). (d) In any SPNE player~2 answers $\hat m_1$ with her
unique best reply $\hat m_2$, so player~1 can secure $\pi_1$ by playing $\hat m_1$; not for every Nash
equilibrium (e.g.\ $(F, (F,F))$ in the example above gives $1 < 2$). (e) Indifferent follower:
$(U,(L,L))$ is an SPNE with payoff $0 < 1$; matching pennies: SPNE payoff $0 <$ mixed Nash payoff
$\frac12$.}
\end{exsol}
